\section{Test und Validierung}

\subsection{Teststrategie und Testumgebung}

Die Validierung erfolgte schrittweise entlang der zentralen
Anwendungsfälle. Zunächst wurden einzelne Komponenten isoliert
überprüft, anschließend das Zusammenspiel von Frontend,
Backend, Workflow Engine und Kommunikationskanälen getestet.

Ein besonderer Fokus lag dabei auf den Übergängen zwischen den
Systemkomponenten, da dort fachliche Fehler oder inkonsistente
Status besonders wahrscheinlich sind.

\subsubsection{Testziele, Testebenen und Auswahl der Prüfschwerpunkte}

\subsubsection{Testwerkzeuge und reproduzierbare Testumgebung}

\subsection{Validierung des Avisierungsservice}

\subsubsection{Domänen- und Anwendungslogik}

\subsubsection{REST-Verträge, Eingabevalidierung und Fehlerantworten}

\subsubsection{Authentifizierung, Sitzungsverwaltung und Mandantenisolation}

\subsubsection{Persistenz- und Abfragemodell}

\subsubsection{Workflow-, Wiederholungs- und Nebenläufigkeitsszenarien}

\subsubsection{Ausgehende Schnittstellen und Rückkanäle}

\subsection{End-to-End-Tests}

Die End-to-End-Tests prüfen die zentralen Benutzerabläufe der
Kundenbestätigung, des Dispositions-Dashboards und der
Lieferverfolgung über Frontend, Backend, E-Mail-Versand und
DISPO-Rückmeldung hinweg. Sie werden mit Playwright im
Chromium-Browser ausgeführt und befinden sich im Verzeichnis
\texttt{frontend/tests}.

\subsubsection{Einordnung der Szenariogruppen}

Die Teststruktur orientiert sich an fachlichen Benutzerabläufen
und nicht an einzelnen Frontend-Komponenten. Die Szenarien sind
in die Bereiche Kundenbestätigung, Dispositions-Dashboard und
Lieferverfolgung gegliedert.

Die Szenarien zur Kundenbestätigung umfassen den Versand und
Inhalt der Bestätigungs-E-Mail, den Zugriff auf die Kundenseite,
die Anzeige der Lieferdaten sowie die Bestätigung oder Ablehnung
eines Liefertermins. Zusätzlich werden ungültige
Bestätigungstoken, optionale Kundenkommentare und die Persistenz
bereits abgegebener Antworten geprüft. Die
Dashboard-Szenarien prüfen den geschützten Zugang, die
Tourenübersicht, die Suche und Filterung, die Auftragsdetails,
die Avisierungshistorie und die Auswahl des Kommunikationskanals
für eine erneute Avisierung. Die Tracking-Szenarien prüfen die
Verfügbarkeit und Aktualisierung der Lieferverfolgung.

\subsubsection{Testumgebung und technische Struktur}

Die Testsuite enthält insgesamt 29 Playwright-Tests. Davon
bilden 28 Tests fachliche End-to-End-Szenarien ab. Ein
zusätzlicher Setup-Test prüft vor der Ausführung die
Erreichbarkeit der benötigten Dienste.

Die fachlichen Tests befinden sich in den Verzeichnissen
\texttt{frontend/tests/confirmation},
\texttt{frontend/tests/dashboard} und
\texttt{frontend/tests/tracking}.

Die Tests zur Kundenbestätigung sind auf folgende Dateien
verteilt:

\begin{itemize}
    \item \texttt{email.spec.ts} für den Versand und Inhalt der
    Bestätigungs-E-Mail
    \item \texttt{page.spec.ts} für die Anzeige der Kundenseite
    und das Verhalten bei ungültigen Bestätigungstoken
    \item \texttt{response.spec.ts} für Bestätigung, Ablehnung,
    Kundenkommentare, DISPO-Rückmeldungen und persistierte
    Abschlusszustände
\end{itemize}

Die Tests des Dispositions-Dashboards befinden sich in folgenden
Dateien:

\begin{itemize}
    \item \texttt{authentication.spec.ts} für den gültigen und
    ungültigen Dashboard-Zugang
    \item \texttt{overview.spec.ts} für die Tourenübersicht,
    Statusanzahlen, Suche, Filterung und Pagination,
    \item \texttt{order-detail.spec.ts} für Auftragsdetails,
    Anfragehistorie, Kundenreaktionen, Resend-Berechtigung und
    Kanalauswahl
    \item \texttt{settings.spec.ts} für das Laden und Testen der
    E-Mail- und SMTP-Einstellungen
\end{itemize}

Die Tests zur Lieferverfolgung befinden sich in
\texttt{frontend/tests/tracking/tracking.spec.ts}. Sie prüfen
die Tracking-Karte, die Fahrzeug- und Zielmarkierung, die
Anzeige der Zieladresse, die Verfügbarkeit vor dem Liefertag
und die Aktualisierung der Fahrerposition.

Wiederverwendbare technische Abläufe befinden sich unter
\texttt{frontend/tests/helpers}. Dazu gehören das Erzeugen
eindeutiger Testaufträge, das Auslesen von E-Mails aus Mailpit,
das Öffnen einer authentifizierten Dashboard-Sitzung, die
Prüfung von DISPO-Rückmeldungen und die Vorbereitung bestätigter
Aufträge für Tracking-Szenarien.

Vor den fachlichen Tests führt Playwright den Setup-Test
\texttt{setup/services.setup.ts} aus. Dieser prüft einmalig die
Erreichbarkeit von Backend, Mailpit und DISPO-Mock. Die
Prüfungen verwenden Polling, damit die Testausführung auf den
Start der benötigten Docker-Dienste warten kann.

Dynamisch erzeugte Testaufträge erhalten mit
\texttt{randomUUID()} eindeutige Auftragsnummern und
E-Mail-Adressen. Die Dashboard-Tests verwenden vorbereitete
\texttt{DEMO}-Seed-Daten. Playwright verändert dabei keine
Datenbankzustände direkt.

\subsubsection{Versand und Bereitstellung der Anfrage}

Zur Prüfung des Versandablaufs wird über die geschützte
DISPO-Schnittstelle eine Bestätigungsanfrage für einen
eindeutigen Testauftrag erzeugt. Die dadurch ausgelöste
Bestätigungs-E-Mail wird über den konfigurierten SMTP-Versand
an Mailpit übermittelt. Dort wartet der Test auf den Eingang
der Nachricht.

Anschließend wird überprüft, ob die E-Mail den Kundennamen, die
Lieferadresse, das Produkt, die Bestellmenge, das
Lieferzeitfenster, den Preis und den personalisierten
Bestätigungslink enthält.

Damit bildet das Szenario den systemübergreifenden Ablauf von
der Anlage eines Avisierungsauftrags bis zur Bereitstellung
eines für den Kunden nutzbaren Bestätigungslinks ab. Bewertet
wird dabei die tatsächlich in Mailpit eingegangene Nachricht.

\subsubsection{Kundenzugriff und Anzeige der Lieferdaten}

Nach dem Öffnen des Bestätigungslinks wird geprüft, ob die
Bestätigungsseite alle für die Kundenentscheidung erforderlichen
Informationen anzeigt. Dazu gehören Kundenname, Lieferadresse,
Produkt, Bestellmenge, Preis sowie Beginn und Ende des
Lieferzeitfensters.

Zusätzlich wird geprüft, wie die Anwendung auf ein unbekanntes
Bestätigungstoken reagiert. Sie zeigt in diesem Fall eine
verständliche Fehlermeldung an und bietet keine Möglichkeit,
den Liefertermin zu bestätigen oder abzulehnen.

\subsubsection{Kundenreaktion und DISPO-Rückmeldung}

Die Bestätigung und die Ablehnung eines Liefertermins werden als
getrennte E2E-Szenarien geprüft.

Nach der Bestätigung muss die Anwendung eine sichtbare
Erfolgsmeldung anzeigen. Zusätzlich wird geprüft, ob der
DISPO-Mock eine Rückmeldung mit der externen Auftragsnummer und
dem Status \texttt{CONFIRMED} erhält.

Ein weiterer Test prüft die Übermittlung einer optionalen
Kundennachricht im Bestätigungsablauf. Der eingegebene Kommentar
muss gemeinsam mit der \texttt{CONFIRMED}-Rückmeldung an den
DISPO-Mock übertragen werden. Da für Bestätigung und Ablehnung
dasselbe Kommentarfeld und derselbe Übertragungsweg verwendet
werden, wird die Kommentarübermittlung nur einmal im
Bestätigungsablauf geprüft.

Nach der Ablehnung muss die Anwendung ebenfalls eine sichtbare
Rückmeldung anzeigen. Der DISPO-Mock muss dabei den Status
\texttt{REJECTED} für die entsprechende externe Auftragsnummer
erhalten.

Die Persistenz beider Abschlusszustände wird durch das erneute
Öffnen des Bestätigungslinks geprüft. Der zuvor erreichte Status
muss weiterhin angezeigt werden, während die Buttons zum
Bestätigen und Ablehnen nicht mehr verfügbar sind. Im
abgelehnten Zustand darf außerdem keine Tracking-Karte
angezeigt werden.

\subsubsection{Dashboard-Zugang}

Für den Dashboard-Zugang fordert der Test mithilfe eines gültigen
API-Schlüssels über die geschützte Backend-Schnittstelle einen
Dashboard-Zugangslink an. Nach dem Öffnen dieses Links im Browser
prüft der Test den Aufbau einer authentifizierten
Dashboard-Sitzung, die Weiterleitung zum Dispositions-Dashboard
und die Anzeige der Tourenübersicht.

Ein weiterer Test ruft die Anwendung mit einem ungültigen
Zugangscode auf. In diesem Fall darf das Dashboard nicht
angezeigt werden. Stattdessen muss die Anwendung eine
verständliche Fehlermeldung ausgeben.

\subsubsection{Tourenübersicht, Suche und Filterung}

In der Tourenübersicht werden die verfügbaren Touren, die
zugehörigen Aufträge und die aggregierten Statusanzahlen
kontrolliert. Dabei werden bestätigte, abgelehnte und noch
unbeantwortete Aufträge berücksichtigt.

Die globale Suche wird in getrennten Szenarien mit einem
Kundennamen, einer Auftragsnummer und einer Lieferadresse
ausgeführt. Dabei wird jeweils geprüft, ob der gesuchte Auftrag
angezeigt und ein nicht passender Auftrag aus der Ergebnisliste
entfernt wird.

Zusätzlich wird die kombinierte Filterung nach Tournummer und
Bestätigungsstatus geprüft. Der Pagination-Test wechselt auf die
zweite Ergebnisseite, kontrolliert den dafür ausgeführten
Backend-Aufruf und prüft die Anzeige der dort enthaltenen
Ergebnisse.

\subsubsection{Auftragsdetails und erneute Avisierung}

Die Auftragsdetailseite zeigt die Stammdaten des Auftrags, den
aktuellen Status und die aktuelle Bestätigungsanfrage. Die
E2E-Tests prüfen unter anderem die Auftragsnummer, den
Kundennamen, die E-Mail-Adresse, die Lieferadresse, das Fahrzeug,
die Bestellmenge und den Kommunikationskanal.

Darüber hinaus wird die Avisierungshistorie geprüft. Historische
Anfragen mit den Zuständen \texttt{REJECTED} und
\texttt{NO\_RESPONSE} sowie gespeicherte Kundenkommentare müssen
auf der Detailseite sichtbar sein.

Die Verfügbarkeit der erneuten Avisierung hängt vom aktuellen
Status ab. Für einen bereits bestätigten Auftrag muss der Button
\glqq Erneut senden\grqq deaktiviert sein. Für einen Auftrag mit
dem Status \texttt{NO\_RESPONSE} muss der Button dagegen
aktiviert sein.

Für eine zulässige erneute Avisierung wird die Auswahl der
Kommunikationskanäle E-Mail, SMS und WhatsApp geprüft. Nach der
Auswahl eines Kanals muss die vorherige Auswahl aufgehoben und
der aktive Kanal über das Attribut \texttt{aria-pressed}
zugänglich gekennzeichnet werden. Der Button zum erneuten Senden
muss dabei aktiviert bleiben.

Der Versand wird im E2E-Test nicht ausgelöst, da die
Dashboard-Tests gemeinsam verwendete \texttt{DEMO}-Seed-Daten
nutzen. Eine erneute Avisierung würde die aktuelle Anfrage und
ihre Historie verändern und dadurch die Wiederholbarkeit der
Tests beeinträchtigen.

Für SMS und WhatsApp wird ausschließlich die Auswahl in der
Benutzeroberfläche geprüft. Ein tatsächlicher Versand über
Twilio erfolgt nicht, um kostenpflichtige Nachrichten an externe
Empfänger und eine Abhängigkeit von der Verfügbarkeit eines
Drittdienstes zu vermeiden.

\subsubsection{E-Mail- und SMTP-Einstellungen}

Die E-Mail-Einstellungen des Dashboards werden in drei
Szenarien geprüft.

Zunächst wird kontrolliert, ob die gespeicherten
SMTP-Einstellungen vom Backend geladen und in der
Benutzeroberfläche dargestellt werden.

Anschließend wird über den Button
\glqq Verbindung testen\grqq{} geprüft, ob das Backend die
gespeicherte SMTP-Konfiguration erfolgreich verwenden kann.
Nach erfolgreichem Abschluss muss eine entsprechende Meldung
angezeigt werden.

Der dritte Test versendet eine SMTP-Test-E-Mail. Neben der
sichtbaren Erfolgsmeldung wird geprüft, ob Mailpit tatsächlich
eine neue Nachricht mit dem erwarteten Betreff erhalten hat.

\subsubsection{Lieferverfolgung und Tracking-Verfügbarkeit}

Für bestätigte Lieferungen am Liefertag wird geprüft, ob die
Live-Tracking-Ansicht verfügbar ist. Die Tracking-Schnittstellen
müssen gültige Ziel- und Fahrerkoordinaten liefern. In der
Benutzeroberfläche müssen die Tracking-Karte sowie die
Fahrzeug- und Zielmarkierung sichtbar sein.

Ein weiteres Szenario prüft, ob die Lieferadresse als
Zieladresse angezeigt wird und die Zielmarkierung auf der Karte
sichtbar ist.

Bei einer bestätigten Lieferung mit einem zukünftigen
Lieferdatum meldet die Tracking-Schnittstelle, dass die
Lieferverfolgung noch nicht verfügbar ist, und liefert keine
Zielkoordinaten. Die Anwendung zeigt daher weder die
Tracking-Karte noch die Fahrzeug- und Zielmarkierung an.
Stattdessen wird der Kunde darauf hingewiesen, dass die
Tracking-Informationen erst am Liefertag verfügbar sind.

Abschließend wird die Aktualisierung der Fahrerposition über den
Button \glqq Aktualisieren\grqq{} geprüft. Dabei müssen neue
Fahrerkoordinaten geladen und der angezeigte Lieferstatus
beziehungsweise die verbleibende Entfernung aktualisiert
werden. Die Karte sowie die Fahrzeug- und Zielmarkierung müssen
weiterhin sichtbar bleiben.
